\section{Introduction}
\label{sec:intro}

Most cryptographic protocols are designed around a single signer. A certificate authority
has one private key; a server authenticating to a client uses one key. This simplicity is
convenient, but it creates a serious problem in settings where the signer is a
high-value target. If that one key is stolen, the attacker can impersonate the signer
forever. If the machine holding the key goes down, signing stops entirely. For a root CA
or a financial institution's signing infrastructure, neither outcome is acceptable.

Threshold signature schemes address this by distributing the signing key across $n$
parties in such a way that any $k$ of them can cooperate to produce a valid signature,
but any coalition smaller than $k$ cannot sign anything, and in fact learns nothing about
the private key. The security guarantee is analogous to Shamir's secret sharing
\cite{Sha79}: the key is shared, not replicated. This is fundamentally different from
simply storing the same key on multiple machines. Threshold security means that even
if up to $k - 1$ machines are fully compromised, the key itself remains secret and the
attacker cannot forge signatures.

\subsection{Shamir Secret Sharing and Its Role in Threshold Signatures}
\label{subsec:sss}

Shamir's Secret Sharing (SSS) \cite{Sha79} is the foundational building block for most
threshold signature schemes. It distributes a secret $s$ among $n$ participants so that
any subset of size at least $t$ (the threshold) can reconstruct $s$, while any subset of
size $t-1$ learns nothing about $s$. The scheme is information‑theoretically secure and
works over a finite field $\mathbb{F}_q$ with $q > n$.

\paragraph{Share generation.}
A dealer chooses a random polynomial $f(x) \in \mathbb{F}_q[x]$ of degree at most $t-1$:
\[
f(x) = a_0 + a_1 x + \cdots + a_{t-1} x^{t-1},
\]
with constant term $a_0 = s$. Coefficients $a_1,\dots,a_{t-1}$ are uniform in $\mathbb{F}_q$.
For each participant $i = 1,\dots,n$, the dealer assigns a distinct non‑zero evaluation
point $x_i \in \mathbb{F}_q^\times$ (public) and computes the share $y_i = f(x_i)$.

\paragraph{Reconstruction via Lagrange interpolation.}
Given any $t$ distinct shares $(x_{i_j}, y_{i_j})$, the polynomial $f$ is uniquely determined.
The secret is $s = f(0) = \sum_{j=1}^{t} y_{i_j} \cdot \ell_j(0)$, where
\[
\ell_j(x) = \prod_{\substack{k=1\\ k \neq j}}^{t} \frac{x - x_{i_k}}{x_{i_j} - x_{i_k}}.
\]

\paragraph{Proof of correctness.}
\textit{Any $t$ distinct shares uniquely determine $s$.} A polynomial of degree $\le t-1$
over a field is uniquely determined by its values at $t$ distinct points. Since $f$ has
degree $\le t-1$ and passes through the given $t$ points, $f(0)=s$ is fixed.

\paragraph{Proof of security (information‑theoretic).}
\textit{Any set of at most $t-1$ shares reveals no information about $s$.}
Fix $t-1$ shares $(x_i, y_i)$. For any candidate secret $s' \in \mathbb{F}_q$, there
exists exactly one polynomial $f'$ of degree $\le t-1$ satisfying $f'(0)=s'$ and
$f'(x_i)=y_i$ (interpolate through the $t$ points $(0,s'), (x_1,y_1),\dots$). Because the
dealer chose the coefficients $a_1,\dots,a_{t-1}$ uniformly, every $s'$ is equally likely
given the $t-1$ shares. Hence the conditional distribution of $s$ is uniform over
$\mathbb{F}_q$; no information is leaked.

\paragraph{From SSS to threshold signatures.}
In a threshold signature scheme, the private signing key $\mathsf{sk}$ is taken as the
secret $s$ and shared using SSS. To sign a message $m$, any $t$ parties each compute a
\textit{partial signature} using their share and a public function (e.g., in Schnorr or
ECDSA, the nonce is also shared). Lagrange interpolation is then applied \emph{in the
exponent} of the cryptographic group to combine the partial signatures into a full
signature. Because $\mathsf{sk}$ is never reconstructed, the scheme remains secure as
long as fewer than $t$ parties are corrupted. The information‑theoretic security of SSS
guarantees that $t-1$ corrupted parties gain no advantage in forging signatures.

\vspace{0.2cm}
\noindent
The topic has direct practical relevance. Certificate authorities using threshold signing
are more resilient to key theft. Cryptocurrency wallets that use threshold signatures
(sometimes called ``multi-sig'') protect funds even if some signers are compromised.
Distributed systems that need to produce signed messages for consensus, as in Byzantine
fault-tolerant protocols, benefit from threshold signatures that can tolerate up to a
third of the parties being malicious. In all of these settings, what is needed is a
scheme that is efficient, non-interactive if possible, and comes with a proof of security
rather than just heuristic arguments.

Constructing such schemes turns out to be harder than it looks. The naive approach of
having all $k$ parties compute partial signatures and combine them runs into the problem
that the combination needs to reconstruct the private key implicitly, without any party
ever computing it explicitly. For RSA-based signatures, this requires doing polynomial
interpolation in a group whose order $\phi(n)$ is secret, which creates technical
difficulties that derailed a decade of prior work. For discrete-log-based signatures, the
inherent randomness in signing forces interaction among the parties to generate shared
nonces. Getting around these obstacles requires non-trivial ideas.

This report covers three papers that together give a solid picture of where threshold
signature research stands.

The first, Shoup's ``Practical Threshold Signatures'' \cite{Shoup00}, is a landmark
result for threshold RSA. The scheme it presents is surprisingly simple, using little more than Shamir secret
sharing and a scaling trick with $l!$ to bypass the $\phi(n)$-interpolation obstacle,
yet it is the first to simultaneously achieve non-interactive signing, compact shares,
and a rigorous proof of security under the RSA assumption. We study this paper in detail in Section~\ref{sec:paper1}, including both
protocols and their security proofs.

The second, Boldyreva's ``Threshold Signatures, Multisignatures and Blind Signatures
Based on the Gap-Diffie-Hellman-Group Signature Scheme'' \cite{Bold03}, takes a
different algebraic route by working in groups where the group order is public (GDH
groups). This lets Lagrange interpolation work directly over $\mathbb{Z}_p$ without any
scaling factor, removes the need for a trusted dealer by using distributed key
generation, and eliminates zero-knowledge proofs from robustness checking. The paper
also extends the same framework to multisignatures and blind signatures. We cover it in
Section~\ref{sec:paper2}.

The third, Komlo and Goldberg's ``FROST: Flexible Round-Optimized Schnorr Threshold
Signatures'' \cite{FROST20}, tackles the round complexity problem for threshold Schnorr
schemes. Prior work required three or more rounds per signing session due to the need
for jointly generated nonces. FROST reduces this to two rounds (or one with a
preprocessing phase) by using non-interactive additive secret sharing combined with
per-participant binding factors that prevent adaptive attacks on the aggregated nonce.
It has been standardised as RFC 9591 and adopted in practice. We cover it in
Section~\ref{sec:frost_review}.

Throughout, we try not just to describe the constructions but to explain the choices made
specifically, why certain assumptions are needed, where the proofs are tight and where
they are not, and what the gaps between theory and practice look like.